% Part: turing-machines 
% Chapter: undecidability 
% Section: verification

\documentclass[../../../include/open-logic-section]{subfiles}

\begin{document}

\olfileid{tur}{und}{ver} 
\olsection{निरूपणाची पडताळणी}

\begin{explain}
आपले निरूपण योग्य आहे हे तपासण्यासाठी दोन
गोष्टी सिद्ध कराव्या लागतील. प्रथम,
आदान~$w$ वर $M$ थांबत असेल, तर
$!T(M, w) \lif !E(M, w)$ वैध आहे
हे दाखवायचे. त्यानंतर उलट दिशाही
दाखवायची: $!T(M, w) \lif !E(M, w)$
वैध असेल, तर आदान~$w$ वर
चालवलेले $M$ खरोखर अखेरीस
थांबते.

या दोन सिद्धतांच्या युक्त्या खूप
वेगळ्या आहेत. पहिल्या फलितासाठी
प्रथम-क्रम तर्कशास्त्रातील
!!a{sentence} वाक्य, म्हणजे
$!T(M, w) \lif !E(M, w)$,
वैध आहे हे दाखवायचे. याचा
सर्वांत सोपा मार्ग म्हणजे
!!a{derivation} निष्पत्ती देणे.
मात्र आपला पुरावा सर्व $M$
आणि $w$ साठी लागू हवा;
म्हणून ज्या एका विशिष्ट
!!{sentence} वाक्याची
!!a{derivation} निष्पत्ती
द्यायची आहे असे नाही,
तर अशी अनंत वाक्ये आहेत.
त्यामुळे $M$ आणि~$w$
कसेही असले आणि आदान~$w$
वर $M$ थांबायला कितीही
पायऱ्या लागत असल्या,
तरी $!T(M, w) \lif !E(M, w)$
याची !!a{derivation} निष्पत्ती
मिळेल, हे गणितीय विगमनाने
सिद्ध करणे हा योग्य मार्ग आहे.

अर्थातच, या विगमनाचा आधार
$M$ थांबणाऱ्या स्थितिवर्णनापर्यंत
पोहोचेपर्यंत केलेल्या पायऱ्यांची
संख्या असेल. विगमनात्मक पुराव्यात
आदान~$w$ वरील $M$ च्या
चालक्रमातील प्रत्येक पायरी~$n$
साठी $!T(M, w) \Entails !C(M, w, n)$
हे सिद्ध करू. येथे $!C(M, w, n)$
हे आदान~$w$ वर चालवलेल्या~$M$
चे $n$ पायऱ्यांनंतरचे स्थितिवर्णन
अचूक सांगते. आता $M$ आदान~$w$
वर, समजा, $n$ पायऱ्यांनंतर
थांबत असेल, तर $!C(M, w, n)$
हे थांबणारे स्थितिवर्णन सांगेल.
असे स्थितिवर्णन सांगणारे
$!C(M, w, n)$ असेल तेव्हा
$!C(M, w, n) \Entails !E(M, w)$
हेही दाखवू. म्हणून $M$
आदान~$w$ वर थांबत असेल,
तर एखाद्या~$n$ साठी
$n$~पायऱ्यांनंतर $M$
थांबणाऱ्या स्थितिवर्णनात असेल.
त्या वेळी $!T(M, w)
\Entails !C(M, w, n)$,
आणि $!C(M, w, n)$
थांबणारे स्थितिवर्णन सांगत
असल्यामुळे $!C(M, w, n) \Entails !E(M, w)$.
त्यामुळे $!T(M, w) \Entails !E(M, w)$,
म्हणजेच $\Entails !T(M, w)
\lif !E(M, w)$ मिळते.

उलट दिशेची युक्ती फार वेगळी
आहे. येथे $\Entails !T(M, w) \lif !E(M, w)$
असे गृहीत धरून $M$ आदान~$w$
वर थांबते हे सिद्ध करायचे.
गृहीतकावरून $!T(M, w) \Entails !E(M,
w)$ मिळते; म्हणजे $!T(M,
w)$ सत्य असलेल्या प्रत्येक
!!{structure} रचनेत
$!E(M, w)$ सत्य आहे.
म्हणून $!T(M, w)$ सत्य
असलेली एक !!a{structure} रचना
~$\Struct{M}$ तयार करू:
तिचे परिभाषाक्षेत्र $\Nat$
असेल, आणि सर्व $\Obj Q_q$
व $\Obj S_\sigma$ यांचे
अर्थनिर्धारण आदान~$w$
वरील~$M$ च्या चालक्रमातील
स्थितिवर्णनांनी ठरेल.
उदाहरणार्थ,
$\Sat{M}{\Obj Q_q(\num{m}, \num{n})}$
तेव्हाच, जेव्हा $M$ हे
आदान~$w$ वर $n$ पायऱ्या
चालवल्यानंतर अवस्था~$q$
मध्ये असून घर~$m$ कडे
पाहत असेल. आता गृहीतकानुसार
$!T(M, w) \Entails !E(M, w)$
आणि रचनेवरून
$\Sat{M}{!T(M,
  w)}$, म्हणून
$\Sat{M}{!E(M, w)}$.
पण $\Sat{M}{!E(M, w)}$
तेव्हाच, जेव्हा एखाद्या
$n \in \Domain{M} = \Nat$
साठी आदान~$w$ वर
चालवलेले $M$ हे
~$n$ पायऱ्यांनंतर
थांबणाऱ्या स्थितिवर्णनात
असेल.
%READERNOTE{OLTUR-013 — गोठवलेल्या इंग्रजी आढाव्यात एका निष्पन्नतेत T(M,w) ऐवजी आवश्यक !T(M,w) मधील ! गाळलेले आहे; पुढील उदाहरणात संपूर्ण बांधणी M साठी असताना अचानक T यंत्र म्हटले आहे. मराठीत दोन्ही ठिकाणी सुसंगत !T(M,w) आणि M वापरले आहेत.}
\end{explain}

\begin{defn} 
$!C(M, w, n)$ हे पुढील
!!{sentence} वाक्य असू द्या:
\[ 
\Obj Q_q(\num{m}, \num{n}) \land \Obj S_{\sigma_0}(\num{0}, \num{n})
\land \dots \land \Obj S_{\sigma_k}(\num{k}, \num{n}) \land
\lforall[x][(\num{k} < x \lif \Obj S_\TMblank(x, \num{n}))]
\] 
येथे वेळ~$n$ ला~$q$
ही $M$ ची अवस्था आहे,
वेळ~$n$ ला $M$ घर~$m$
कडे पाहत आहे, वेळ~$n$
ला $0 \le i \le k$
साठी घर~$i$ मध्ये
चिन्ह~$\sigma_i$ आहे,
आणि $k$ हे वेळ~$0$
ला फितीवरील सर्वांत
उजवे रिकामे नसलेले
घर किंवा $n$ पायऱ्यांपर्यंत
शीर्षाने भेट दिलेले
सर्वांत उजवे घर,
यांपैकी जे अधिक
उजवीकडे असेल ते आहे.
\end{defn}

\begin{lem}\ollabel{lem:halt-config-implies-halt}
आदान~$w$ वर चालवलेले $M$
हे $n$ पायऱ्यांनंतर
थांबणाऱ्या स्थितिवर्णनात
असेल, तर $!C(M, w, n) \Entails !E(M, w)$.
\end{lem}

\begin{proof}
आदान~$w$ वर $M$ हे $n$
पायऱ्यांनंतर थांबते असे
समजा. मग अशी अवस्था~$q$,
घर~$m$ आणि चिन्ह~$\sigma$
असतील की
\begin{enumerate} 
\item $n$ पायऱ्यांनंतर $M$
  हे~$q$ अवस्थेत असून,
  ज्या घर~$m$ वर~$\sigma$
  आहे त्याकडे पाहत असेल.
\item संक्रमण फलन $\delta(q, \sigma)$
  अपरिभाषित असेल.
\end{enumerate}
$!C(M, w, n)$ हे या
स्थितिवर्णनाचे वर्णन आहे;
त्यात $\Obj Q_{q}(\num{m}, \num{n})$
आणि $\Obj
S_{\sigma}(\num{m}, \num{n})$
ही संयुक्त विधाने असतील.
त्यांवरून एकत्रितपणे
$!E(M,
w)$ मिळते:
\[
\lexists[x][\lexists[y][(\bigvee_{\tuple{q, \sigma} \in
      X}(\Obj Q_q(x, y) \land \Obj S_\sigma(x, y)))]]
\]
कारण $\tuple{q,\sigma} \in X$
असल्याने
$\Obj Q_{q}(\num{m}, \num{n}) \land \Obj S_{\sigma}(\num{m},
\num{n}) \Entails \bigvee_{\tuple{q, \sigma} \in X} (\Obj Q_q(\num{m},
\num{n}) \land \Obj S_{\sigma}(\num{m}, \num{n}))$.
%READERNOTE{OLTUR-012 — गोठवलेल्या इंग्रजी पुराव्यात आधी थांबणारी जोडी (q,σ) निवडल्यानंतर या निष्पन्नतेत अचानक (q',σ') येते, आणि एका S-विधेयापूर्वी Obj चिन्हही गाळले आहे. येथे आधी निवडलेलीच (q,σ) जोडी व भाषेतील तेच S-विधेय वापरले आहे; अस्तित्वविधानाचा निष्कर्ष बदलत नाही.}
\end{proof}

\begin{explain}
म्हणून $M$ आदान $w$
वर थांबत असेल, तर
एखाद्या~$n$ साठी
$!C(M,
w, n) \Entails !E(M,w)$.
आता चालक्रम पोहोचतो अशा कोणत्याही वेळ~$n$
साठी $!T(M, w) \Entails !C(M, w, n)$
हे दाखवू.
\end{explain}

\begin{lem}
\ollabel{lem:config}
प्रत्येक $n$ साठी,
$M$ चा चालक्रम $n$ पायऱ्यांनंतरच्या
स्थितिवर्णनापर्यंत पोहोचत असेल, तर
$!T(M, w)
\Entails !C(M, w, n)$.
\end{lem}
%READERNOTE{OLTUR-011 — गोठवलेल्या इंग्रजी लेम्मात यंत्र n पायऱ्यांनंतर थांबलेले नसावे अशी अट आहे, पण पुढील पुरावा n व्या पायरीवर मिळणाऱ्या थांबणाऱ्या स्थितिवर्णनासाठीही तेच लेम्मा वापरतो. येथे चालक्रम n-पायऱ्यांच्या स्थितिवर्णनापर्यंत पोहोचणे ही अचूक अट घेतली आहे; ते स्थितिवर्णन थांबणारे असू शकते.}

\begin{proof}
विगमनाची आरंभीची पायरी:
$n = 0$ असेल, तर
$!C(M, w, 0)$ मधील
संयुक्त विधाने $!T(M, w)$
मध्येही आहेत; म्हणून
ती त्यावरून निष्पन्न होतात.

विगमनाचे गृहीतक:
$n$-व्या पायरीपूर्वी $M$
थांबलेले नसेल, तर
$!T(M,w) \Entails !C(M, w, n)$.
आता $!C(M, w, n)$
थांबणारे स्थितिवर्णन सांगत
नसेल, तर $!T(M, w) \Entails
!C(M, w, n+1)$ हे
दाखवायचे आहे.

$n \ge 0$ असे समजा
आणि आदान $w$ वर
सुरू केलेले $M$ हे
$n$ पायऱ्यांनंतर
अवस्था~$q$ मध्ये असून
घर~$m$ कडे पाहत
असेल. $n$ पायऱ्यांनंतर
$M$ थांबत नसल्यामुळे
$M$ च्या कार्यक्रमात
पुढील तीनपैकी एका
रूपाची सूचना असली पाहिजे:

\begin{enumerate} 
\item \ollabel{right} $\delta(q, \sigma) = \tuple{q', \sigma', \TMright}$

\item \ollabel{left} $\delta(q, \sigma) = \tuple{q', \sigma', \TMleft}$

\item \ollabel{stay} $\delta(q, \sigma) = \tuple{q', \sigma', \TMstay}$
\end{enumerate}

या तीनही बाबींचा आता
एकेक करून विचार करू.

\begin{enumerate} 
\item \olref{right} या रूपाची सूचना
  आहे असे समजा. अशा वेळी
  \olref[rep]{defn:tm-descr}\olref[rep]{rep-right}
  यांनुसार
\begin{align*} 
& \lforall[x][\lforall[y][((\Obj Q_{q}(x,
  y) \land \Obj S_\sigma(x, y)) \lif {}]]\\
& \qquad (\Obj Q_{q'}(x',
  y') \land \Obj S_{\sigma'}(x, y') \land !A(x, y)))  
\intertext{हे $!T(M,w)$ चे एक संयुक्त विधान आहे.
  त्यावरून पुढील !!{sentence} वाक्य
  मिळते (सार्वत्रिक विधानात~$x$ ऐवजी
  $\num{m}$ आणि~$y$ ऐवजी
  $\num{n}$ बसवून):}
& (\Obj Q_{q}(\num{m}, \num{n}) \land \Obj S_{\sigma}(\num{m},
\num{n})) \lif {}\\
& \qquad (\Obj Q_{q'}(\num{m}', \num{n}') \land
\Obj S_{\sigma'}(\num{m}, \num{n}') \land !A(\num{m}, \num{n})).
\intertext{विगमनाच्या गृहीतकानुसार
  $!T(M, w) \Entails !C(M, w, n)$,
  म्हणजे}
& \Obj Q_q(\num{m}, \num{n}) \land \Obj S_{\sigma_0}(\num{0}, \num{n})
\land \dots \land \Obj S_{\sigma_k}(\num{k}, \num{n}) \land \\
& \qquad\lforall[x][(\num{k} < x \lif \Obj S_\TMblank(x, \num{n}))]\\
\intertext{$n$ पायऱ्यांनंतर फितीवरील घर~$m$ मध्ये
  चिन्ह~$\sigma$ असल्याने, त्याच्याशी संबंधित
  संयुक्त विधान~$\Obj S_\sigma(\num{m}, \num{n})$
  आहे; म्हणून यावरून पुढील मिळते:}
& \Obj Q_{q}(\num{m}, \num{n}) \land \Obj S_{\sigma}(\num{m},
   \num{n})
\intertext{आता आपल्याला पुढील मिळते:}
& \Obj Q_{q'}(\num{m}', \num{n}') \land \Obj S_{\sigma'}(\num{m},
  \num{n}') \land {}\\
& \qquad\Obj S_{\sigma_0^{+}}(\num{0}, \num{n}') \land \dots \land
  \Obj S_{\sigma_k^{+}}(\num{k}, \num{n}') \land {}\\
& \qquad \lforall[x][(\num{k} < x \lif \Obj S_\TMblank(x, \num{n}'))]
\end{align*}
येथे वरची अधिक खूण पुढील स्थितीतील चिन्ह
दर्शवते: लिहिलेल्या घरात नव्या सूचनेतील चिन्ह,
आणि इतर घरांत आधीची चिन्हे आहेत.
याचे कारण असे: पहिली ओळ आधीच्या
सोपाधिक विधानाच्या परिणामी भागावरून
निगमन नियमाने थेट मिळते.
मधल्या ओळीतील प्रत्येक संयुक्त
विधान---ज्यात
$S_{\sigma_m}(\num{m},\num{n}')$
वगळले आहे---हे~$!C(M, w, n)$
मधील संबंधित संयुक्त विधान
आणि $!A(\num{m},
\num{n})$ यांवरून मिळते.

$m < k$ असेल, तर
$!T(M,w) \Proves \num{m} < \num {k}$
(\olref[rep]{prop:mlessk}),
आणि~$<$ च्या संक्रमणशीलतेवरून
$\lforall[x][(\num{k} < x \lif \num{m} < x)]$
मिळते. $m = k$ असेल,
तर $\lforall[x][(\num{k} < x \lif \num{m} < x)]$
केवळ तर्कनियमांवरून मिळते.
म्हणून शेवटची ओळ
$!C(M, w,
n)$ मधील संबंधित संयुक्त विधान,
$\lforall[x][(\num{k} < x \lif \num{m} < x)]$,
आणि $!A(\num{m},
\num{n})$ यांवरून मिळते.
$m<k$ असेल, तर हेच
$!C(M, w, n+1)$ आहे.

आता $m=k$ असे समजा.
या वेळी $n+1$ पायऱ्यांनंतर
शीर्षाने घर~$k+1$ लाही
भेट दिलेली असेल; तेच
आता भेट दिलेले सर्वांत
उजवे घर आहे. म्हणून
$!C(M, w, n+1)$ मध्ये
$\Obj
S_\TMblank(\num{k}',\num{n}')$
हे नवे संयुक्त विधान असेल,
आणि शेवटचे संयुक्त विधान
$\lforall[x][(\num{k}' < x \lif \Obj S_\TMblank(x, \num{n}'))]$
असेल. ही दोन्ही !!{sentence}s
देखील निष्पन्न होतात
हे तपासावे लागेल.

आपल्याकडे आधीच
$\lforall[x][(\num{k} < x \lif \Obj S_\TMblank(x,
  \num{n}'))]$
आहे. विशेषतः त्यावरून
$\num{k} < \num{k}' \lif
\Obj S_\TMblank(\num{k}', \num{n}')$
मिळते. स्वयंसिद्धक
$\lforall[x][x <
  x']$ यावरून
$\num{k} < \num{k}'$
मिळते. सोपाधिक
निगमनाने
$\Obj
S_\TMblank(\num{k}',\num{n}')$
मिळते.

शिवाय,
$!T(M,w) \Proves \num{k} < \num{k}'$
असल्यामुळे~$<$ च्या
संक्रमणशीलतेच्या
स्वयंसिद्धकावरून
$\lforall[x][(\num{k}' < x \lif \Obj
  S_\TMblank(x, \num{n}'))]$
मिळते. (याची पडताळणी
सरावासाठी सोडतो.)

\item \olref{left} या रूपाची
सूचना आहे असे समजा.
मग \olref[rep]{defn:tm-descr}\olref[rep]{rep-left}
यांनुसार
\begin{align*} 
& \lforall[x][\lforall[y][((\Obj Q_{q}(x', y) \land \Obj
    S_{\sigma}(x', y)) \lif {}]]\\
& \qquad (\Obj Q_{q'}(x, y') \land \Obj
  S_{\sigma'}(x', y') \land !A(x', y))) \land {}\\
& \lforall[y][((\Obj Q_{q_i}(\num{0}, y) \land \Obj S_{\sigma}(\num{0},
    y)) \lif {}]\\
& \qquad (\Obj Q_{q_j}(\num{0}, y') \land \Obj S_{\sigma'}(\num{0}, y')
  \land !A(\num{0}, y)))
\intertext{हे $!T(M,w)$ चे एक
  संयुक्त विधान आहे. $m>0$
  असेल, तर $l = m - 1$
  असे घ्या (म्हणजे
  $m = l+1$). वरील
  !!{sentence} वाक्याच्या
  पहिल्या संयुक्त विधानावरून
  पुढील मिळते:}
& (\Obj Q_{q}(\num{l}', \num{n}) \land \Obj S_{\sigma}(\num{l}', \num{n}))
\lif {} \\
& \qquad (\Obj Q_{q'}(\num{l}, \num{n}') \land \Obj S_{\sigma'}(\num{l}',
\num{n}') \land !A(\num{l}', \num{n}))
\intertext{अन्यथा $l = m = 0$
  असे घ्या आणि दुसऱ्या
  संयुक्त विधानावरून निष्पन्न
  होणारे पुढील !!{sentence}
  वाक्य पाहा:}
& ((\Obj Q_{q_i}(\num{0}, \num{n}) \land \Obj S_{\sigma}(\num{0}, \num{n})) \lif {}\\
& \qquad   (\Obj Q_{q_j}(\num{0}, \num{n}') \land \Obj S_{\sigma'}(\num{0}, \num{n}') \land
!A(\num{0}, \num{n}))) 
\intertext{शून्यावरील शाखेतील पहिली अवस्था
  आधीची, तर दुसरी पुढची अवस्था आहे.
  दोन्हीपैकी कोणत्याही
  वाक्यावरून पुढील मिळते:}
&  \Obj Q_{q'}(\num{l}, \num{n}') \land \Obj S_{\sigma'}(\num{m},
  \num{n}') \land {}\\
&\qquad  \Obj S_{\sigma_0^{+}}(\num{0}, \num{n}')\land \dots \land
  \Obj S_{\sigma_k^{+}}(\num{k}, \num{n}')  \land {} \\
& \qquad  \lforall[x][(\num{k} < x
  \lif \Obj S_\TMblank(x, \num{n}'))]
\end{align*}
पहिल्या शाखेत स्थिरता-नियम लिहिले गेलेले
मूळ घर वगळतो; डावीकडे गेलेले नवे शीर्षस्थान नाही.
येथेही वरची अधिक खूण पुढील स्थितीतील
चिन्ह दर्शवते; लिहिलेल्या घरात नवे चिन्ह
आणि इतर घरांत आधीची चिन्हे आहेत.
%READERNOTE{OLTUR-014 — गोठवलेल्या इंग्रजी पडताळणीत डावीकडील संक्रमणाच्या पहिल्या शाखेत A(x,y) आणि त्याच्या उदाहरणात A(l,n) पुन्हा आले आहे. OLTUR-009 नुसार लिहिलेले घर x'=l'=m वगळले पाहिजे. म्हणून येथे A(x',y) आणि A(l',n) वापरले आहेत; घर शून्यवरील शाखा बदललेली नाही.}
%READERNOTE{OLTUR-015 — गोठवलेल्या इंग्रजी उजवीकडील व डावीकडील पुढील-स्थिती प्रदर्शनांत प्रत्येक घरासाठी जुने σ_i चिन्हच पुन्हा दाखवले आहे, जरी आधीचे वाक्य लिहिलेल्या घरात σ' येते असे म्हणते आणि पुढील स्पष्टीकरण जुने त्या घराचे विधान वगळते. येथे वरची + खूण पुढील-स्थिती चिन्हांसाठी आहे: लिहिलेल्या घरात नवे चिन्ह आणि इतरत्र जुने चिन्ह. त्यामुळे प्रदर्शनाचे वाचन स्थितिवर्णनाच्या व्याख्येशी जुळते.}
आधीप्रमाणेच. (पहिल्या बाबतीत
$\num{l}' \ident \num{l+1}
\ident \num{m}$ आणि
दुसऱ्या बाबतीत
$\num{l} \ident \num{0}$
हे लक्षात घ्या.) पण
हेच $!C(M, w, n+1)$
आहे.

\item \olref{stay} ही बाब
सरावासाठी सोडतो.
\end{enumerate}
म्हणून चालक्रम पोहोचतो अशा कोणत्याही~$n$
साठी $!T(M, w) \Entails !C(M, w, n)$
हे दाखवले.
\end{proof}

\begin{prob}
\olref[tur][und][ver]{lem:config}
च्या पुराव्यातील
~\olref[tur][und][ver]{stay}
ही बाब पूर्ण करा.
\end{prob}

\begin{prob}
$\Obj S_{\sigma_i}(\num{i}, \num{n})$
आणि $!A(m, n)$ यांवरून
$\Obj S_{\sigma_i}(\num{i}, \num{n}')$
ची !!a{derivation} निष्पत्ती द्या
($i \neq
m$, म्हणजे $i < m$
किंवा $m < i$,
असे गृहीत धरून).
\end{prob}

\begin{prob}
$\lforall[x][(\num{k} < x \lif \Obj
  S_\TMblank(x, \num{n}'))]$,
$\lforall[x][x < x']$,
आणि
$\lforall[x][\lforall[y][\lforall[z][ ((x < y \land y < z) \lif x <
      z)]]]$
यांवरून
$\lforall[x][(\num{k}' < x \lif \Obj
  S_\TMblank(x, \num{n}'))]$
ची !!a{derivation}
निष्पत्ती द्या.
\end{prob}


\begin{lem}
\ollabel{lem:valid-if-halt}
$M$ आदान~$w$ वर
थांबत असेल, तर
$!T(M, w) \lif
!E(M, w)$ वैध आहे.
\end{lem}

\begin{proof}
\olref{lem:config} नुसार,
चालक्रम पोहोचतो अशा कोणत्याही वेळ~$n$
ला~$M$ चे वेळ~$n$ वरील स्थितिवर्णन
सांगणारे~$!C(M, w, n)$
हे~$!T(M, w)$ वरून
निष्पन्न होते. समजा,
$M$ हे $k$ पायऱ्यांनंतर
थांबते. त्या वेळी
ते कोणत्यातरी~$m \in \Nat$
साठी घर~$m$ कडे
पाहत असेल. मग
$!C(M, w, k)$ हे~$M$
चे थांबणारे स्थितिवर्णन
सांगते. म्हणजे त्यात
$\Obj Q_q(\num{m}, \num{k})$
आणि $\Obj
S_\sigma(\num{m}, \num{k})$
ही दोन्ही संयुक्त विधाने
असतील व $\delta(q,\sigma)$
अपरिभाषित असेल. म्हणून
\olref{lem:halt-config-implies-halt}
नुसार
$!C(M, w, k) \Entails !E(M,
w)$. पण
$!T(M, w) \Entails !C(M, w, k)$
असल्यामुळे
$!T(M, w)
\Entails !E(M, w)$;
म्हणून $!T(M, w) \lif !E(M, w)$
वैध आहे.
\end{proof}


\begin{explain} 
आपला दावा पूर्णपणे
पडताळण्यासाठी उलटी
दिशाही स्थापित करावी
लागेल: $!T(M, w) \lif !E(M, w)$
वैध असेल, तर
आदान~$w$ वर सुरू
केलेले $M$ खरोखर
थांबते.
\end{explain}

\begin{lem}
\ollabel{lem:halt-if-valid}
$\Entails !T(M, w) \lif !E(M, w)$
असेल, तर $M$
आदान~$w$ वर थांबते.
\end{lem}

\begin{proof}
परिभाषाक्षेत्र~$\Nat$ असलेली
$\Lang L_M$-!!{structure} रचना
~$\Struct M$ घ्या. ती
$\Obj 0$ चा अर्थ~$0$,
$\prime$ चा अर्थ उत्तरवर्ती
फलन आणि $<$ चा अर्थ
लहान-असणे संबंध असा
लावते. $\Obj Q_q$
आणि~$\Obj S_\sigma$
ही विधेये पुढीलप्रमाणे
अर्थनिर्धारित करतात:
\begin{align*}
  \Assign{\Obj Q_q}{M} & =
\Setabs{\tuple{m, n}}{\begin{array}{ll}\text{आदान~$w$ वर सुरू होऊन $n$ पायऱ्यांनंतर,}\\ \text{$M$ हे अवस्था $q$
  मध्ये असून घर~$m$ कडे पाहते}\end{array}} \\
\Assign{\Obj S_\sigma}{M} & = \Setabs{\tuple{m, n}}{\begin{array}{ll}
\text{आदान~$w$ वर सुरू होऊन $n$ पायऱ्यांनंतर,}\\ \text{$M$ च्या घर~$m$ मध्ये
  चिन्ह~$\sigma$ आहे}\end{array}}
\end{align*}
म्हणजे, आदान~$w$ वर
सुरू केलेले $M$ प्रत्यक्षात
प्रत्येक पायरीला काय करते
हे वर्णन करणारी !!{structure}
रचना~$\Struct{M}$
आपण बांधतो. स्पष्टपणे,
$\Sat{M}{!T(M, w)}$.
$\Entails !T(M, w) \lif !E(M, w)$
असेल, तर
$\Sat{M}{!E(M, w)}$
देखील; म्हणजे
\[
\Sat{M}{\lexists[x][\lexists[y][(\bigvee_{\tuple{q, \sigma} \in
      X}(\Obj Q_q(x, y) \land \Obj S_\sigma(x, y)))]]}.
\]
$\Domain{M} = \Nat$
असल्यामुळे अशी
$m$, $n \in \Nat$
असतील की काही~$q$
आणि~$\sigma$ साठी
$\delta(q, \sigma)$
अपरिभाषित असून
$\Sat{M}{\Obj Q_q(\num{m}, \num{n}) \land \Obj S_\sigma(\num{m},
\num{n})}$
असेल. ~$\Struct M$
च्या व्याख्येनुसार,
याचा अर्थ आदान~$w$
वर सुरू केलेले $M$
हे~$n$ पायऱ्यांनंतर
अवस्था~$q$ मध्ये असून
चिन्ह~$\sigma$ वाचत आहे,
आणि संक्रमण फलन
अपरिभाषित आहे;
म्हणजे~$M$ थांबले आहे.
\end{proof}

\end{document}
